% Template for PLoS
% Version 3.5 March 2018
%
% % % % % % % % % % % % % % % % % % % % % %
%
% -- IMPORTANT NOTE
%
% This template contains comments intended
% to minimize problems and delays during our production
% process. Please follow the template instructions
% whenever possible.
%
% % % % % % % % % % % % % % % % % % % % % % %
%
% Once your paper is accepted for publication,
% PLEASE REMOVE ALL TRACKED CHANGES in this file
% and leave only the final text of your manuscript.
% PLOS recommends the use of latexdiff to track changes during review, as this will help to maintain a clean tex file.
% Visit https://www.ctan.org/pkg/latexdiff?lang=en for info or contact us at latex@plos.org.
%
%
% There are no restrictions on package use within the LaTeX files except that
% no packages listed in the template may be deleted.
%
% Please do not include colors or graphics in the text.
%
% The manuscript LaTeX source should be contained within a single file (do not use \input, \externaldocument, or similar commands).
%
% % % % % % % % % % % % % % % % % % % % % % %
%
% -- FIGURES AND TABLES
%
% Please include tables/figure captions directly after the paragraph where they are first cited in the text.
%
% DO NOT INCLUDE GRAPHICS IN YOUR MANUSCRIPT
% - Figures should be uploaded separately from your manuscript file.
% - Figures generated using LaTeX should be extracted and removed from the PDF before submission.
% - Figures containing multiple panels/subfigures must be combined into one image file before submission.
% For figure citations, please use "Fig" instead of "Figure".
% See http://journals.plos.org/plosone/s/figures for PLOS figure guidelines.
%
% Tables should be cell-based and may not contain:
% - spacing/line breaks within cells to alter layout or alignment
% - do not nest tabular environments (no tabular environments within tabular environments)
% - no graphics or colored text (cell background color/shading OK)
% See http://journals.plos.org/plosone/s/tables for table guidelines.
%
% For tables that exceed the width of the text column, use the adjustwidth environment as illustrated in the example table in text below.
%
% % % % % % % % % % % % % % % % % % % % % % % %
%
% -- EQUATIONS, MATH SYMBOLS, SUBSCRIPTS, AND SUPERSCRIPTS
%
% IMPORTANT
% Below are a few tips to help format your equations and other special characters according to our specifications. For more tips to help reduce the possibility of formatting errors during conversion, please see our LaTeX guidelines at http://journals.plos.org/plosone/s/latex
%
% For inline equations, please be sure to include all portions of an equation in the math environment.
%
% Do not include text that is not math in the math environment.
%
% Please add line breaks to long display equations when possible in order to fit size of the column.
%
% For inline equations, please do not include punctuation (commas, etc) within the math environment unless this is part of the equation.
%
% When adding superscript or subscripts outside of brackets/braces, please group using {}.
%
% Do not use \cal for caligraphic font.  Instead, use \mathcal{}
%
% % % % % % % % % % % % % % % % % % % % % % % %
%
% Please contact latex@plos.org with any questions.
%
% % % % % % % % % % % % % % % % % % % % % % % %

\documentclass[10pt,letterpaper]{article}
\usepackage[top=0.85in,left=2.75in,footskip=0.75in]{geometry}

% amsmath and amssymb packages, useful for mathematical formulas and symbols
\usepackage{amsmath,amssymb}

% Use adjustwidth environment to exceed column width (see example table in text)
\usepackage{changepage}

% Use Unicode characters when possible
\usepackage[utf8x]{inputenc}

% textcomp package and marvosym package for additional characters
\usepackage{textcomp,marvosym}

% cite package, to clean up citations in the main text. Do not remove.
% \usepackage{cite}

% Use nameref to cite supporting information files (see Supporting Information section for more info)
\usepackage{nameref,hyperref}

% line numbers
\usepackage[right]{lineno}

% ligatures disabled
\usepackage{microtype}
\DisableLigatures[f]{encoding = *, family = * }

% color can be used to apply background shading to table cells only
\usepackage[table]{xcolor}

% array package and thick rules for tables
\usepackage{array}

% create "+" rule type for thick vertical lines
\newcolumntype{+}{!{\vrule width 2pt}}

% create \thickcline for thick horizontal lines of variable length
\newlength\savedwidth
\newcommand\thickcline[1]{%
  \noalign{\global\savedwidth\arrayrulewidth\global\arrayrulewidth 2pt}%
  \cline{#1}%
  \noalign{\vskip\arrayrulewidth}%
  \noalign{\global\arrayrulewidth\savedwidth}%
}

% \thickhline command for thick horizontal lines that span the table
\newcommand\thickhline{\noalign{\global\savedwidth\arrayrulewidth\global\arrayrulewidth 2pt}%
\hline
\noalign{\global\arrayrulewidth\savedwidth}}


% Remove comment for double spacing
%\usepackage{setspace}
%\doublespacing

% Text layout
\raggedright
\setlength{\parindent}{0.5cm}
\textwidth 5.25in
\textheight 8.75in

% Bold the 'Figure #' in the caption and separate it from the title/caption with a period
% Captions will be left justified
\usepackage[aboveskip=1pt,labelfont=bf,labelsep=period,justification=raggedright,singlelinecheck=off]{caption}
\renewcommand{\figurename}{Fig}

% Use the PLoS provided BiBTeX style
% \bibliographystyle{plos2015}

% Remove brackets from numbering in List of References
\makeatletter
\renewcommand{\@biblabel}[1]{\quad#1.}
\makeatother



% Header and Footer with logo
\usepackage{lastpage,fancyhdr,graphicx}
\usepackage{epstopdf}
%\pagestyle{myheadings}
\pagestyle{fancy}
\fancyhf{}
%\setlength{\headheight}{27.023pt}
%\lhead{\includegraphics[width=2.0in]{PLOS-submission.eps}}
\rfoot{\thepage/\pageref{LastPage}}
\renewcommand{\headrulewidth}{0pt}
\renewcommand{\footrule}{\hrule height 2pt \vspace{2mm}}
\fancyheadoffset[L]{2.25in}
\fancyfootoffset[L]{2.25in}
\lfoot{\today}

%% Include all macros below

\newcommand{\lorem}{{\bf LOREM}}
\newcommand{\ipsum}{{\bf IPSUM}}






\usepackage{forarray}
\usepackage{xstring}
\newcommand{\getIndex}[2]{
  \ForEach{,}{\IfEq{#1}{\thislevelitem}{\number\thislevelcount\ExitForEach}{}}{#2}
}

\setcounter{secnumdepth}{0}

\newcommand{\getAff}[1]{
  \getIndex{#1}{Department of Hospital Pathology, Seoul St.~Mary's Hospital}
}

\providecommand{\tightlist}{%
  \setlength{\itemsep}{0pt}\setlength{\parskip}{0pt}}

\begin{document}
\vspace*{0.2in}

% Title must be 250 characters or less.
\begin{flushleft}
{\Large
\textbf\newline{Prediction of \emph{PIK3CA} mutations from cancer gene expression data} % Please use "sentence case" for title and headings (capitalize only the first word in a title (or heading), the first word in a subtitle (or subheading), and any proper nouns).
}
\newline
% Insert author names, affiliations and corresponding author email (do not include titles, positions, or degrees).
\\
Jun Kang\textsuperscript{\getAff{Department of Hospital Pathology, Seoul St.~Mary's Hospital}},
Ahwon Lee\textsuperscript{\getAff{Department of Hospital Pathology, Seoul St.~Mary's Hospital}},
Youn Soo Lee\textsuperscript{\getAff{Department of Hospital Pathology, Seoul St.~Mary's Hospital}}\textsuperscript{*}\\
\bigskip
\textbf{\getAff{Department of Hospital Pathology, Seoul St.~Mary's Hospital}}Department of Hospital Pathology, Seoul St.~Mary's Hospital, College of
Medicine, The Catholic University of Korea, Seoul, South Korea\\
\bigskip
* Corresponding author: lys9908@catholic.ac.kr\\
\end{flushleft}
% Please keep the abstract below 300 words
\section*{Abstract}
Breast cancers with \emph{PIK3CA} mutations can be treated with
\emph{PIK3CA} inhibitors in hormone receptor-positive HER2 negative
subtypes. We applied a supervised elastic net penalized logistic
regression model to predict \emph{PIK3CA} mutations from gene expression
data. This regression approach was applied to predict modeling using the
TCGA pan-cancer dataset. Approximately 10,000 cases were available for
\emph{PIK3CA} mutation and mRNA expression data. In 10-fold
cross-validation, the model with \(\lambda\) = 0.01 and \(\alpha\) = 1.0
(ridge regression) showed the best performance, in terms of area under
the receiver operating characteristic (AUROC). The final model was
developed with selected hyper-parameters using the entire training set.
The training set AUROC was 0.93, and the test set AUROC was 0.84. The
area under the precision-recall (AUPR) of the training set was 0.66, and
the test set AUPR was 0.39. Cancer types were the most important
predictors. Both \emph{insulin like growth factor 1 receptor}
(\emph{IGF1R}) and the \emph{phosphatase and tensin homolog}
(\emph{PTEN}) were the most significant genes in gene expression
predictors. Our study suggests that predicting genomic alterations using
gene expression data is possible, with good outcomes.

% Please keep the Author Summary between 150 and 200 words
% Use first person. PLOS ONE authors please skip this step.
% Author Summary not valid for PLOS ONE submissions.

\linenumbers

% Use "Eq" instead of "Equation" for equation citations.
\hypertarget{introduction}{%
\section{Introduction}\label{introduction}}

Targeted therapy has become a standard treatment for many cancer
patients, however the approach requires a test for a specific cancer
genomic alteration, to treat patients. Several direct genomic alteration
tests have been developed and proven for their clinical utility to treat
patients {[}1{]}.

Machine learning approaches can be applied to detect genomic
alterations. Machine learning algorithms can build prediction models
from a large number of predictors, such as radiomic features {[}3{]},
pathology image {[}4{]} or gene expression data {[}5{]}. Because most
direct genomic tests are more specific and sensitive than predictive
models, machine learning approaches may have limited roles in clinical
practice, however, machine learning approaches are ideal when direct
tests are unavailable or fail.

RAS pathway activation predictions have been performed using gene
expression data {[}5{]}. Authors used data from The Cancer Genome Atlas
(TCGA), with a supervised elastic net penalized logistic regression
classifier, with stochastic gradient descent. Their model performance
was 84\% with an area under the receiver operating characteristic
(AUROC) curve, and 63\% with an area under the precision-recall (AUPR)
curve. Importantly, these authors suggested their approach could be
applied to other genomic alterations.

Breast cancer having \emph{PIK3CA} mutations can be treated using PIK3CA
inhibitors, in hormone receptor-positive HER2 negative subtypes {[}7{]}.
The \emph{PIK3CA} mutation is the second most common driver mutation
after \emph{TP53}, and is most frequently detected in endometrial
carcinoma (45\%), followed by breast invasive carcinoma (24\%), cervical
squamous cell carcinoma, endo-cervical adenocarcinoma (20\%) and colon
adenocarcinoma (16\%) {[}8{]}.

\emph{PIK3CA} encodes the p110\(\alpha\) catalytic subunit of
phosphatidylinositol 3′-kinase (PI3K). PI3K is a protein kinase that
phosphorylates phosphatidylinositol 4,5-biphosphate
(PIP\textsubscript{2}) to generate phosphatidylinositol
3,4,5-triphosphate (PIP\textsubscript{3}). The phosphatase and tensin
homolog (\emph{PTEN}) converts PIP\textsubscript{2} to
PIP\textsubscript{3} in contrast to PI3K. PIP\textsubscript{3} is a
second messenger that activates protein kinase B (AKT), which is a
serine/threonine-specific protein kinase. AKT inhibits apoptosis and
promotes cell proliferation {[}9{]}.

We applied a supervised elastic net penalized logistic regression model
to predict \emph{PIK3CA} mutations. We wanted to ascertain whether this
prediction model approach could be applied not only to RAS pathway
activation, but also to \emph{PIK3CA} mutation predictions. The purpose
of this study is to investigate the \emph{PIK3CA} mutation prediction
performance of machine learning models.

\hypertarget{materials-and-methods}{%
\section{Materials and Methods}\label{materials-and-methods}}

\hypertarget{dataset}{%
\subsection{Dataset}\label{dataset}}

We used the TCGA pan-cancer dataset. TCGA archives the following; exome
sequencing, gene expression, DNA methylation, protein expression, and
clinical data from \textgreater{} 10,000 cancer samples across 33 common
cancer types. The TCGA dataset is publically available. \emph{PIK3CA}
mutation data was extracted using cgdsr rpackage {[}10{]}. Gene
expression data was downloaded from the National Cancer Institute
(NCI)'s Genomic Data Commons (GDC) website. This archives data for TCGA
(https://gdc.cancer.gov/about-data/publications/pancanatlas). Gene
expression in the TCGA pan-cancer dataset is batch-corrected with
normalization.

The target variable was \emph{PIK3CA} mutation status. \emph{PIK3CA}
status was considered positive when the case had the following
\emph{PIK3CA} variants (C420R, E542K, E545A, E545D, E545G, E545K, Q546E,
Q546R, H1047L, H1047R, H1047Y) which were the target variants of the
Therascreen \emph{PIK3CA} RGQ PCR Kit, Qiagen, Hilden, Germany. This kit
was approved as a companion diagnostics test to treat with PIK3CA
inhibitor by the United States Food and Drug Administration.

\hypertarget{modeling-process}{%
\subsection{Modeling process}\label{modeling-process}}

To narrow down potential predictors, genes with a large median absolute
deviation (\textgreater{} third-quartiles) were selected. Thirty three
cancer type dummy variables were included in predictor variables. We
split three-quarters of the dataset into the training set and one
quarter into the test set. Yeo-Johnson transformation was performed to
correct skewness. Centering and scaling were also performed. All
preprocessing was performed using the recipe r package {[}11{]}.
Penalized logistic regression was applied to prediction modeling.
Ten-fold cross-validation with target variable stratification was
performed over the hyper-parameter grid: \(\lambda\) \{\(10^{-5}\),
\(10^{-4}\),\(10^{-3}\),\(10^{-2}\),\(10^{-1}\), \(10^{0}\)\},
\(\alpha\) \{0.0, 0.25, 0.5, 0.75, 1.0\}. Lambda (\(\lambda\)) is a
penalty scaling parameter and alpha (\(\alpha\)) is a mixing parameter
of penalty function
(\((1-\alpha)/2 \lVert\beta\rVert_2^2+\alpha\lVert \beta \rVert_1\))
{[}12{]}.

\hypertarget{assessing-model-performance}{%
\subsection{Assessing model
performance}\label{assessing-model-performance}}

Model performance was evaluated using AUROC and AUPR curve approaches.
The AUPR approach is more informative than AUROC for imbalanced datasets
{[}13{]}. The modeling process and assessing model performance were
performed with the tidymodels rpackage {[}14{]}.

\hypertarget{results}{%
\section{Results}\label{results}}

\hypertarget{dataset-summary}{%
\subsection{Dataset summary}\label{dataset-summary}}

10,845 cases were available for both \emph{PIK3CA} mutation and mRNA
expression data. 5,128 out of 20,502 genes were included in the modeling
process, after filtering for median absolute deviation, as described in
the modeling process method. The prevalence rate for \emph{PIK3CA}
mutation was 0.11 in all cases. The \emph{PIK3CA} mutation prevalence
rate in each cancer type varied. The median prevalence rate of
\emph{PIK3CA} mutation for each cancer type was 0.03 (range 0--0.33)
(Figure 1).

\hypertarget{selecting-model-and-performance-estimation}{%
\subsection{Selecting model and performance
estimation}\label{selecting-model-and-performance-estimation}}

For 10-fold cross-validation, the model with \(\lambda\) = 0.01 and
\(\alpha\) = 1.0 (ridge regression) showed the best performance in terms
of AUROC (S1 Figure). The final model was trained with the selected
hyper-parameters with the entire training set. The training set AUROC
was 0.93 and the test set AUROC was 0.84. The AUPR of the training set
was 0.66 and the test set AUPR was 0.39 (Figure 2A).

\hypertarget{performance-of-each-cancer-type}{%
\subsection{Performance of each cancer
type}\label{performance-of-each-cancer-type}}

Because \emph{PIK3CA} mutation prevalence varied across cancer types,
the performance of each cancer type was investigated. The AUROC and AUPR
were positively correlated between the training sets and test sets in
cancer type sub-analysis (Figure 2B). The AUPR was high in cancer types
with high \emph{PIK3CA} mutation rates such as colon, breast and uterus
cancer types. The AUROC did not correlate with \emph{PIK3CA} mutation
rates of each cancer type (Figure 2C).

\hypertarget{important-predictors}{%
\subsection{Important predictors}\label{important-predictors}}

The top 30 important predictors are shown (Figure 3). The coefficient is
the parameter of the predictor which represents the effect of the
predictor on prediction. \emph{Insulin like growth factor 1 Receptor}
(\emph{IGF1R}) mRNA expression was the strongest negative predictor, and
\emph{PTEN} was the strongest positive predictor. Both \emph{IGF1R} and
\emph{PTEN} are key players in the tyrosine kinase pathway {[}9,15{]}.
The cancer types were important predictors. Some cancer types including
uterine carcinosarcoma (UCS), bladder urothelial carcinoma (BLCA),
pancreatic adenocarcinoma (PAAD), lymphoid neoplasm diffuse large B-cell
lymphoma (DLBC) were the strongest predictors.

\hypertarget{discussion}{%
\section{Discussion}\label{discussion}}

Our model showed good performance in predicting \emph{PIK3CA} mutations
in various cancer types. Our data suggested that the supervised elastic
net penalized logistic regression model could be applied not only to the
RAS activation pathway, but also to other genomic alterations. Both the
RAS activation pathway and \emph{PIK3CA} mutations are key, common
cancer genomic alterations. Because they exert significant effect on
gene expression in cancer cells, prediction from gene expression data
can be good. However, the supervised elastic net penalized logistic
regression model cannot be generalized or applied to other genomic
alterations which have have a weak effect on gene expression.

Prediction modeling from the TCGA pan-cancer dataset can be limiting in
terms of data preprocessing. The gene expression data is processed by
between-sample normalization to remove batch effects. If the model has
been trained from between-sample normalization, a new sample cannot be
exactly processed with normalization which was done on trainset. A model
based on gene expression from the TCGA pan-cancer dataset has limitation
in terms of data preprocessing. It is necessary to develop a processing
method that is independent of a dataset, to apply gene expression data
to the prediction model.

Our \emph{PIK3CA} prediction model was similar to the RAS activation
prediction model in terms of AUROC (0.84). However the AUPR of our model
was lower than the RAS activation model (0.39 versus 0.63). The reason
for our lower AUPR may be explained by an imbalanced dataset that has
the low prevalence rate of \emph{PIK3CA} mutations {[}5{]}. The model
for RAS activation trained with cancer types with more than 0.05
prevalence of RAS activation to avoid imbalance classification problem.
We included all cancer types in our modeling process. The lower
prevalence rate of target variables meant our dataset had a lower AUPR
baseline. In the sub-analysis performance of each cancer type, the
cancer types with higher \emph{PIK3CA} mutation rates showed better
AUPRs.

Our model included cancer types as predictors, and they were stronger
predictors than gene expression. The varying prevalence of PIC3CA
mutations across cancer types may be a reason for the strong predictive
power of cancer types.

Some significant gene expression predictors were closely related to the
PTEN-PI3K pathway. \emph{PTEN} and \emph{IGFR1R} were the strongest gene
expression predictors, which has negative and positive predictive
powers. \emph{IGF1R} is a tyrosine kinase receptor that activates PI3K
{[}15{]}, and \emph{PTEN} is an important regulator of
PIP\textsubscript{3} by dephosphorylating PIP\textsubscript{3} {[}9{]}.

Several studies have attempted to predict genomic alterations from gene
expression data {[}6,16{]}. A study investigated \emph{PIK3CA} mutation
predictions using gene-expression signatures which is a sum of the
average of the logarithmic gene expression. The model showed good
performance AUROC 0.71 in an independent test set {[}6,16{]}. Another
study predicted copy number alterations with gene expression, using a
multinomial logistic regression model with least absolute shrinkage and
selection operator (LASSO) parameters {[}17{]}. The prediction of the
1p/19q codeletion was very good, with an AUROC of 0.997, and gene-level
predictions were good, with an AUROC of 0.75 {[}17{]}. A logistic
regression model was used for \emph{MYCN Proto-Oncogene, BHLH
Transcription Factor} (\emph{MYCN}) gene amplification in neuroblastoma
{[}18{]}.

The clinical utility of PIK3CA mutation prediction from mRNA expression
is unclear because most direct genomic tests are more specific and
sensitive than predictive models. Our prediction model is not an
application that is immediately applicable to a cancer patient for
detection of PIK3CA mutation. It is not known how it will be used, but
finding out the mutation prediction performance using gene expression
data could play a role in advancing machine learning to be helpful in
patient treatment.

Our study suggested that the prediction of genomic alterations using
gene expression data was possible, with good performance. However,
improved performances are required for clinical tests, and the
standardization of generation processing of gene expression data is also
needed.

\hypertarget{figure-legends}{%
\section{Figure legends}\label{figure-legends}}

\begin{itemize}
\item
  Figure 1. Prevalence rate of \emph{PIK3CA} mutations across cancer
  types. Cancer type abbreviations are explained in the S1 Appendix.
\item
  Figure 2. Summary of modeling results. (A) Left: receiver operating
  characteristic (ROC) curve. Right: precision-recall (PR) curve of
  training set and test set. The horizontal green line is the
  \emph{PIK3CA} mutation rate (0.11) (B) Correlation between training
  set and test set of the area under the receiver operating
  characteristic curve (AUROC), and the area under the precision-recall
  curve (AUPR) among cancer types. The gray band is the 95\% confidence
  interval. Abbreviations are explained in the S1 Appendix. (C)
  Correlations between the \emph{PIK3CA} mutation rate of the AUROC, and
  the AUPR.
\item
  Figure 3. Coefficient model. (A) Top 30 high mRNA coefficients. (B)
  Cancer type coefficients. Cancer types abbreviations are explained in
  the S1 Appendix.
\end{itemize}

\hypertarget{supporting-information}{%
\section{Supporting information}\label{supporting-information}}

\begin{itemize}
\tightlist
\item
  S1 Appendix. Abbreviations of cancer type.\\
\item
  S1 Figure. Hyperparameter tuning and performance assessment in 10-fold
  cross-validation resampling. The x-axis is a penalty scaling
  parameter: \(\lambda\) \{\(10^{-5}\),
  \(10^{-4}\),\(10^{-3}\),\(10^{-2}\),\(10^{-1}\), \(10^{0}\)\}, color
  is mixture hyperparameter of penalty function: \(\alpha\) \{0.0, 0.25,
  0.5, 0.75, 1.0\}. y-axis is estimates of area under the receiver
  operating characteristic (AUROC) using 10-fold cross-validation
  resampling.
\end{itemize}

\hypertarget{references}{%
\section{References}\label{references}}

\hypertarget{refs}{}
\leavevmode\hypertarget{ref-healthNucleicAcidBased2020}{}%
1. Health C for D and R. Nucleic Acid Based Tests. FDA. FDA;

\leavevmode\hypertarget{ref-sahnanePyrosequencingEGFRMutation2013}{}%
2. Sahnane N, Gueli R, Tibiletti M, Bernasconi B, Stefanoli M, Franzi F,
et al. Pyrosequencing for EGFR Mutation Detection: Diagnostic Accuracy
and Clinical Implications. Diagnostic Molecular Pathology. 2013;22:
196--203.
doi:\href{https://doi.org/10.1097/PDM.0b013e3182893f55}{10.1097/PDM.0b013e3182893f55}

\leavevmode\hypertarget{ref-dercleIdentificationNonSmall2020}{}%
3. Dercle L, Fronheiser M, Lu L, Du S, Hayes W, Leung DK, et al.
Identification of NonSmall Cell Lung Cancer Sensitive to Systemic Cancer
Therapies Using Radiomics. Clin Cancer Res. American Association for
Cancer Research; 2020;26: 2151--2162.
doi:\href{https://doi.org/10.1158/1078-0432.CCR-19-2942}{10.1158/1078-0432.CCR-19-2942}

\leavevmode\hypertarget{ref-coudrayClassificationMutationPrediction2018}{}%
4. Coudray N, Ocampo PS, Sakellaropoulos T, Narula N, Snuderl M, Fenyö
D, et al. Classification and mutation prediction from nonSmall cell lung
cancer histopathology images using deep learning. Nature Medicine.
Nature Publishing Group; 2018;24: 1559--1567.
doi:\href{https://doi.org/10.1038/s41591-018-0177-5}{10.1038/s41591-018-0177-5}

\leavevmode\hypertarget{ref-wayMachineLearningDetects2018}{}%
5. Way GP, Sanchez-Vega F, La K, Armenia J, Chatila WK, Luna A, et al.
Machine Learning Detects Pan-cancer Ras Pathway Activation in The Cancer
Genome Atlas. Cell Reports. Elsevier; 2018;23: 172--180.e3.
doi:\href{https://doi.org/10.1016/j.celrep.2018.03.046}{10.1016/j.celrep.2018.03.046}

\leavevmode\hypertarget{ref-loiPIK3CAMutationsAssociated2010}{}%
6. Loi S, Haibe-Kains B, Majjaj S, Lallemand F, Durbecq V, Larsimont D,
et al. PIK3CA mutations associated with gene signature of low mTORC1
signaling and better outcomes in estrogen receptorPositive breast
cancer. PNAS. National Academy of Sciences; 2010;107: 10208--10213.
doi:\href{https://doi.org/10.1073/pnas.0907011107}{10.1073/pnas.0907011107}

\leavevmode\hypertarget{ref-andre_2019_alpelisib_nengljmeda}{}%
7. André F, Ciruelos E, Rubovszky G, Campone M, Loibl S, Rugo HS, et al.
Alpelisib for PIK3CA-Mutated, Hormone ReceptorPositive Advanced Breast
Cancer. New England Journal of Medicine. Massachusetts Medical Society;
2019;380: 1929--1940.
doi:\href{https://doi.org/10.1056/NEJMoa1813904}{10.1056/NEJMoa1813904}

\leavevmode\hypertarget{ref-baileyComprehensiveCharacterizationCancer2018}{}%
8. Bailey MH, Tokheim C, Porta-Pardo E, Sengupta S, Bertrand.
Comprehensive Characterization of Cancer Driver Genes and Mutations.
Cell. 2018;173: 371--385.e18.
doi:\href{https://doi.org/10.1016/j.cell.2018.02.060}{10.1016/j.cell.2018.02.060}

\leavevmode\hypertarget{ref-cantleyPhosphoinositide3KinasePathway2002}{}%
9. Cantley LC. The Phosphoinositide 3-Kinase Pathway. Science. American
Association for the Advancement of Science; 2002;296: 1655--1657.
doi:\href{https://doi.org/10.1126/science.296.5573.1655}{10.1126/science.296.5573.1655}

\leavevmode\hypertarget{ref-jacobsenCgdsrRbasedAPI2019}{}%
10. Jacobsen A, Luna A. Cgdsr: R-based API for accessing the MSKCC
cancer genomics data server (CGDS). 2019.

\leavevmode\hypertarget{ref-kuhnRecipesPreprocessingTools2020}{}%
11. Kuhn M, Wickham H. Recipes: Preprocessing tools to create design
matrices. 2020.

\leavevmode\hypertarget{ref-friedmanRegularizationPathsGeneralized2010}{}%
12. Friedman J, Hastie T, Tibshirani R. Regularization Paths for
Generalized Linear Models via Coordinate Descent. J Stat Softw. 2010;33:
1--22.

\leavevmode\hypertarget{ref-saitoPrecisionRecallPlotMore2015}{}%
13. Saito T, Rehmsmeier M. The Precision-Recall Plot Is More Informative
than the ROC Plot When Evaluating Binary Classifiers on Imbalanced
Datasets. PLOS ONE. Public Library of Science; 10: e0118432.
doi:\href{https://doi.org/10.1371/journal.pone.0118432}{10.1371/journal.pone.0118432}

\leavevmode\hypertarget{ref-kuhnTidymodelsEasilyInstall2020}{}%
14. Kuhn M, Wickham H. Tidymodels: Easily install and load the
'tidymodels' packages. 2020.

\leavevmode\hypertarget{ref-leroithInsulinlikeGrowthFactor2003}{}%
15. LeRoith D, Roberts CT. The insulin-like growth factor system and
cancer. Cancer Letters. 2003;195: 127--137.
doi:\href{https://doi.org/10.1016/S0304-3835(03)00159-9}{10.1016/S0304-3835(03)00159-9}

\leavevmode\hypertarget{ref-cizkovaGeneExpressionProfiling2010}{}%
16. Cizkova M, Cizeron-Clairac G, Vacher S, Susini A, Andrieu C,
Lidereau R, et al. Gene Expression Profiling Reveals New Aspects of
PIK3CA Mutation in ERalpha-Positive Breast Cancer: Major Implication of
the Wnt Signaling Pathway. PLOS ONE. Public Library of Science; 5:
e15647.
doi:\href{https://doi.org/10.1371/journal.pone.0015647}{10.1371/journal.pone.0015647}

\leavevmode\hypertarget{ref-muCNAPEMachineLearning2019}{}%
17. Mu Q, Wang J. CNAPE: A Machine Learning Method for Copy Number
Alteration Prediction from Gene Expression. IEEE/ACM Transactions on
Computational Biology and Bioinformatics. 2019; 1--1.
doi:\href{https://doi.org/10.1109/TCBB.2019.2944827}{10.1109/TCBB.2019.2944827}

\leavevmode\hypertarget{ref-heGeneSignaturesAssociated2020}{}%
18. He X, Qin C, Zhao Y, Zou L, Zhao H, Cheng C. Gene signatures
associated with genomic aberrations predict prognosis in neuroblastoma.
Cancer Communications. 2020;40: 105--118.
doi:\href{https://doi.org/10.1002/cac2.12016}{10.1002/cac2.12016}

\nolinenumbers


\end{document}

